\subsubsection*{Connective Bases}

% BEGIN GENERATED ARTIFACT: def:connective-basis-propositional-logic
\begin{definitionbox}{Definition (Connective Basis)}
\begin{definition}[Connective Basis]
\label{def:connective-basis-propositional-logic}
A \emph{connective basis} is a specified set \(B\) of propositional connectives
allowed as primitive connectives for building formulas.
\end{definition}
\end{definitionbox}

\begin{remark*}[Standard quantified statement]
\[
B\text{ is a connective basis}
\Longleftrightarrow
B\subseteq\{\neg,\land,\lor,\to,\leftrightarrow,\ldots\}
\]
and formulas over \(B\) are built using only connectives from \(B\).
\end{remark*}

\begin{remark*}[Predicate reading]
\(\operatorname{ConnectiveBasis}(B)\) means that \(B\) is the chosen primitive
connective vocabulary for a propositional language or fragment.
\end{remark*}

\begin{remark*}[Interpretation]
A connective basis controls which connectives are primitive, not which truth
functions are expressible. Different bases may have the same expressive power.
\end{remark*}

\begin{dependencies}
\begin{itemize}
  \item \hyperref[def:logical-connective]{Logical Connective}
  \item \hyperref[def:propositional-language]{Propositional Language}
\end{itemize}
\end{dependencies}
% END GENERATED ARTIFACT: def:connective-basis-propositional-logic

% BEGIN GENERATED ARTIFACT: def:definable-connective-propositional-logic
\begin{definitionbox}{Definition (Definable Connective)}
\begin{definition}[Definable Connective]
\label{def:definable-connective-propositional-logic}
Let \(B\) be a connective basis. A connective \(\circ\) is \emph{definable from}
\(B\) if there is a formula built using only connectives from \(B\) whose truth
function agrees with the truth function of \(\circ\).
\end{definition}
\end{definitionbox}

\begin{remark*}[Standard quantified statement]
\[
\circ\text{ is definable from }B
\Longleftrightarrow
\exists\theta\in\WFF_B\;\bigl(f_{\theta}=f_{\circ}\bigr),
\]
where \(\WFF_B\) denotes the formulas built using only connectives from \(B\).
\end{remark*}

\begin{remark*}[Predicate reading]
\(\operatorname{DefinableFrom}(\circ,B)\) means that \(\circ\) can be treated as
an abbreviation for a formula pattern using only connectives from \(B\).
\end{remark*}

\begin{remark*}[Interpretation]
Definability separates notation from expressive power. A connective can be
omitted from the primitive syntax without losing its use, provided it is
recoverable as a defined connective.
\end{remark*}

\begin{dependencies}
\begin{itemize}
  \item \hyperref[def:connective-basis-propositional-logic]{Connective Basis}
  \item \hyperref[def:truth-function-induced-by-formula]{Truth Function Induced by a Formula}
  \item \hyperref[def:boolean-function-propositional-logic]{Boolean Function}
\end{itemize}
\end{dependencies}
% END GENERATED ARTIFACT: def:definable-connective-propositional-logic

% BEGIN GENERATED ARTIFACT: def:functionally-complete-connective-set-propositional-logic
\begin{definitionbox}{Definition (Functionally Complete Connective Set)}
\begin{definition}[Functionally Complete Connective Set]
\label{def:functionally-complete-connective-set-propositional-logic}
A connective basis \(B\) is \emph{functionally complete} if every Boolean
function \(f:\mathbb{B}^n\to\mathbb{B}\), for every \(n\ge 1\), is represented
by some formula using only connectives from \(B\).
\end{definition}
\end{definitionbox}

\begin{remark*}[Standard quantified statement]
\[
\operatorname{FunctionallyComplete}(B)
\Longleftrightarrow
(\forall n\ge 1)(\forall f:\mathbb{B}^n\to\mathbb{B})(\exists\varphi\in\WFF_B)
\operatorname{Represents}(\varphi,f;P_1,\ldots,P_n).
\]
\end{remark*}

\begin{remark*}[Predicate reading]
\(\operatorname{FunctionallyComplete}(B)\) means that formulas over \(B\) can
represent every finite Boolean input-output table.
\end{remark*}

\begin{remark*}[Interpretation]
Functional completeness is the central expressive-power property for connective
sets. It says that the basis can express all possible truth-functional behavior,
not merely the connectives originally listed in the language.
\end{remark*}

\begin{dependencies}
\begin{itemize}
  \item \hyperref[def:connective-basis-propositional-logic]{Connective Basis}
  \item \hyperref[def:boolean-function-representation-propositional-logic]{Representation of a Boolean Function}
  \item \hyperref[def:boolean-function-propositional-logic]{Boolean Function}
\end{itemize}
\end{dependencies}
% END GENERATED ARTIFACT: def:functionally-complete-connective-set-propositional-logic

% BEGIN GENERATED ARTIFACT: def:expressively-equivalent-connective-sets-propositional-logic
\begin{definitionbox}{Definition (Expressively Equivalent Connective Sets)}
\begin{definition}[Expressively Equivalent Connective Sets]
\label{def:expressively-equivalent-connective-sets-propositional-logic}
Two connective bases \(B\) and \(C\) are \emph{expressively equivalent} if every
connective in \(B\) is definable from \(C\), and every connective in \(C\) is
definable from \(B\).
\end{definition}
\end{definitionbox}

\begin{remark*}[Standard quantified statement]
\[
B\equiv_{\mathrm{expr}} C
\Longleftrightarrow
\bigl((\forall\circ\in B)\operatorname{DefinableFrom}(\circ,C)\bigr)
\land
\bigl((\forall\star\in C)\operatorname{DefinableFrom}(\star,B)\bigr).
\]
\end{remark*}

\begin{remark*}[Predicate reading]
\(B\equiv_{\mathrm{expr}} C\) means that the two bases can define exactly the
same connective behavior.
\end{remark*}

\begin{remark*}[Interpretation]
Expressive equivalence lets a proof replace one primitive vocabulary with
another. This is why a language can be simplified syntactically without changing
which truth functions can be expressed.
\end{remark*}

\begin{dependencies}
\begin{itemize}
  \item \hyperref[def:connective-basis-propositional-logic]{Connective Basis}
  \item \hyperref[def:definable-connective-propositional-logic]{Definable Connective}
\end{itemize}
\end{dependencies}
% END GENERATED ARTIFACT: def:expressively-equivalent-connective-sets-propositional-logic
